% answers_short.tex —— 2007 年·普通高等学校招生考试·福建卷·文科 · 简短答案速查
% 由 tools/make_short_answers.py 自 answers.tex 的 \gksolhead/\gkanswer 提取，
% 并与 2007/2007_数学真题_福建文_答案要点.md 逐题比对后生成，勿手改（改源 MD 后重跑）。

\begin{gknotice}
\textbf{使用说明}：本节只给最终答案便于快速核对，完整推导见解析版。答案为编写者独立推导，\textbf{非官方标准答案}。
\end{gknotice}

\gkpart{选择题}
\begin{center}\begin{tabular}{ccccccccccccc}
\toprule
题号 & 1 & 2 & 3 & 4 & 5 & 6 & 7 & 8 & 9 & 10 & 11 & 12 \\
\midrule
答案 & C & C & D & A & A & B & D & B & D & B & B & C \\
\bottomrule
\end{tabular}\end{center}

\gkpart{填空题}
\begin{itemize}[leftmargin=2.6em]
  \item[\textbf{13.}] $15$。
  \item[\textbf{14.}] $[-5,7]$。
  \item[\textbf{15.}] $\dfrac12$。
  \item[\textbf{16.}] 例如：三角形的相似关系；整数的同余关系（模 $n$ 同余）。
\end{itemize}

\gkpart{解答题}
\begin{itemize}[leftmargin=2.6em]
  \item[\textbf{17.}] （1）$C=135^\circ$；（2）$BC=\sqrt2$。
  \item[\textbf{18.}] （1）$0.063$；（2）$0.88$；（3）$0.3024$。
  \item[\textbf{19.}] （1）证明见解析；（2）$\arccos\dfrac{\sqrt6}4$。
  \item[\textbf{20.}] （1）$h(t)=-t^3+t-1$；（2）$m>1$。
  \item[\textbf{21.}] （1）$a_1=1$，$a_n=2\cdot3^{n-2}$（$n\geqslant2$）；（2）$T_n=\dfrac{(2n-1)3^n+3}6$。
  \item[\textbf{22.}] （1）$y^2=4x$；（2）① $0$；② $16$。
\end{itemize}
